% Generated draft: requires mathematical and source review.
\documentclass[11pt]{article}
\usepackage[margin=25mm]{geometry}
\usepackage{fontspec}
\setmainfont{Latin Modern Roman}
\usepackage{amsmath,amssymb,mathtools,unicode-math}
\setmathfont{Latin Modern Math}
\usepackage{enumitem,needspace,xurl,hyperref}
\setlength{\parindent}{0pt}
\setlength{\parskip}{5pt}
\setlength{\emergencystretch}{3em}
\newcommand{\N}{\mathbb{N}}
\newcommand{\Z}{\mathbb{Z}}
\newcommand{\Q}{\mathbb{Q}}
\newcommand{\R}{\mathbb{R}}
\newcommand{\C}{\mathbb{C}}
\newcommand{\F}{\mathbb{F}}
\newcommand{\A}{\mathbb{A}}
\newcommand{\PP}{\mathbb{P}}
\newcommand{\E}{\mathbb{E}}
\newcommand{\eps}{\varepsilon}
\newcommand{\dd}{\,\mathrm{d}}
\begin{document}
\section*{UnsolvedMath ID: 20000966}
\textbf{Source problem code:} catalogue entry
\subsection{Definitions and mathematical statement}
The source summarizes Spencer Backman's proposal for chip-firing on a $k$-uniform hypergraph with a cyclic ordering on each edge and chip values among $k$th roots of unity. It reports a positive-semidefinite complex Laplacian and invariance under the total vertex ordering, asks for its kernel and notions of nonnegativity, reduced divisors and recurrence, and refers to an unspecified statement of interest. Andrea Sportiello suggests replacing roots by general vectors summing to zero. The suggested argument sector for ``nonnegative'' is typographically incomplete, and the source does not supply a legal-firing rule. Its reported quotations are explicitly summaries, not verbatim statements.

For the attached aid's incidence-Gram model let $k\ge2$, $\zeta=e^{2\pi i/k}$, and let an edge have cyclic order $(v_0,\ldots,v_{k-1})$. For $r\in S\subseteq\{1,\ldots,k-1\}$ define $b_{e,r}\in\C^V$ by $b_{e,r}(v_j)=\zeta^{rj}$ and zero off that edge. With positive weights set
\[
L=\sum_{e\in E}\sum_{r\in S}w_r b_{e,r}b_{e,r}^*.
\]
Every mode has coordinate sum zero, so $L\mathbf1=0$, and $L$ is Hermitian positive semidefinite. Changing the cyclic starting point only changes each vector by a unit phase; changing the global vertex enumeration permutes coordinates. This algebraic invariance does not yet establish order independence of an unspecified legal-firing process.

If the bipartite incidence graph is a tree with $m=|E|$ hyperedges, then $|V|=1+(k-1)m$. Writing $s=|S|$, the aid asserts
\[
\operatorname{rank}_{\C}L=sm,\qquad\dim_{\C}\ker L=1+(k-1-s)m.
\]
For cyclotomic modules choose, for example, positive integer weights, so entries lie in $\mathcal O=\Z[\zeta]$, and write $\mathcal O_0^V=\{x:\sum_vx_v=0\}$. The degree-zero firing quotient
\[
\mathcal O_0^V/L\mathcal O^V
\]
has $\mathcal O$-rank $(k-1-s)m$, meaning rank after tensoring with $\Q(\zeta)$. Thus a single nontrivial mode on a nontrivial incidence tree gives an infinite quotient whenever $k\ge3$.

Using every nontrivial mode with positive weights, the aid states $\ker L=\C\mathbf1$ for every incidence-connected hypergraph. With cyclotomic-integral weights, each one-sink reduced matrix then gives a finite module $\mathcal O^{V\setminus\{o\}}/L_{\mathrm{red}}\mathcal O^{V\setminus\{o\}}$. Unequal positive weights on conjugate modes can retain cyclic orientation information that equal full-mode weights lose. Full mode count is therefore necessary for finite graph-style quotients uniformly over incidence trees and sufficient for these finite algebraic modules on connected hypergraphs. It does not by itself answer the original positivity, reduced-divisor or recurrent-configuration questions.
\subsection{Short English statement}
Generalize chip-firing using cyclic roots-of-unity edge data. The supplied aid shows why one Fourier mode leaves extra kernel and an infinite tree quotient, while all nontrivial modes give the constant kernel and finite sink modules under integral weights. Legal dynamics and positivity remain unspecified.
\subsection{Sources}
\begin{itemize}
\item Local UnsolvedMath registry, permanent ID $20000966$, source code catalogue entry; canonical \texttt{aim-combinatorics-notes.json}, \texttt{notes[90]}.
\item Sam Hopkins, recorder, \emph{Problems from the AIM Chip-Firing Workshop}, first problem, following ``Generalizations of chip-firing and the critical group,'' July 8--12, 2013: \url{https://aimath.org/pastworkshops/chipfiringproblems.pdf}. The source explicitly warns that quoted remarks summarize rather than directly quote participants. It attributes the proposal to Spencer Backman and the zero-sum-vector suggestion to Andrea Sportiello. No PDF inspection is claimed here.
\item Backman's statement of interest on the workshop website is mentioned without a precise reference. Attached research attempt 1: partial result, run and registry status partially solved, rigor proved, novelty confidence low, classification PARTIAL-PROGRESS. The machine-generated aid is not peer reviewed; independently verify before citation.
\end{itemize}
\end{document}
